\documentclass{abntex2}
\usepackage[utf8]{inputenc}
\usepackage[T1]{fontenc}
\usepackage{geometry}
\usepackage{graphicx}
\usepackage{float}
\usepackage{hyperref}
\usepackage[english,indonesian]{babel}

\title{Optimasi Pemrosesan Template LaTeX Menggunakan Pendekatan Berbasis AI}
\author{Academic Support Agent}
\date{\today}

\begin{document}
\maketitle

\begin{abstract}

\end{abstract}

\chapter{PENDAHULUAN}


\chapter{TINJAUAN PUSTAKA}


% Original content (commented out):
% %******************************************************************************************
% INSTITUTO FEDERAL DE EDUCAÇÃO, CIÊNCIA E TECNOLOGIA DE PERNAMBUCO

% MODELO DE RELATÓRIO DE PRÁTICAS DO IFPE CAMPUS PESQUEIRA

% DESENVOLVIDO POR CHRISTIAN ARAUJO, EGRESSO DO IFPE CAMPUS PESQUEIRA
% Autor: Christian Araujo
% email: cas@discente.ifpe.edu.br
%******************************************************************************************

% Caso precise de packages adicionais para o seu texto, coloque no arquivo "packages.tex" que está na pasta "Estrutura".

% As figuras vão na pasta "Figuras".

% Os arquivos da pasta "Texto" são os arquivos que você deve editar para adequar seu arquivo da maneira que necessita.

% ---- Pasta texto ----

% O arquivo "dados.tex" é o editável com as informações que vão para a capa e folha de rosto.

% O arquivo "referencias.bib" é o editável que armazena as informações sobre as referências utilizadas no texto.

% O arquivo "texto.tex" é o editável onde você escreve seu texto do relatório.
%******************************************************************************************

\documentclass[
    12pt,
    openright,
    oneside,
    a4paper,
    chapter=TITLE,
    sumario=tradicional,
    english,
    brazil
]{abntex2}

\input{Estrutura/Modelo}

\begin{document}
% Retira espaço extra obsoleto entre as frases.
\frenchspacing
\pagenumbering{gobble}

\capa

\folhaderosto

\input{Texto/texto}

\end{document}


\end{document}